\textbf{1. Kiến thức cần nhớ}
\begin{itemize}
\item \textbf{Các bước giải bài toán bằng cách lập hệ phương trình}:
\begin{enumerate}
\item \textbf{Bước 1. Lập hệ phương trình}:
\begin{itemize}
\item Chọn hai ẩn số biểu thị hai đại lượng chưa biết và đặt \textbf{đơn vị, điều kiện thích hợp} cho các ẩn.
\item Biểu diễn các đại lượng chưa biết khác theo các ẩn và các số liệu đã biết.
\item Dựa vào hai mối quan hệ độc lập trong đề bài để lập hai phương trình, từ đó có hệ phương trình.
\end{itemize}
\item \textbf{Bước 2. Giải hệ phương trình}: Sử dụng phương pháp cộng đại số hoặc phương pháp thế.
\item \textbf{Bước 3. Trả lời}: Đối chiếu nghiệm của hệ với điều kiện của ẩn và kết luận bài toán.
\end{enumerate}
\end{itemize}

\textbf{2. Phương pháp giải các dạng}

%% DANG 9D1?3-1
\textbf{Dạng 1. Toán chuyển động}
\begin{enumerate}
\item Công thức cốt lõi: $\boxed{s = v \cdot t} \iff v = \dfrac{s}{t} \iff t = \dfrac{s}{v}$.
\item Chuyển động trên dòng nước:
$$v_{\text{xuôi dòng}} = v_{\text{thực}} + v_{\text{dòng nước}};\quad v_{\text{ngược dòng}} = v_{\text{thực}} - v_{\text{dòng nước}}.$$
\item Hai xe đi ngược chiều gặp nhau: $s_1 + s_2 = s_{\text{toàn đoạn}}$.
\item Hai xe đi cùng chiều đuổi kịp nhau: $s_1 - s_2 = \text{khoảng cách ban đầu}$.
\end{enumerate}
Lưu ý: Luôn đổi đồng nhất đơn vị (giờ - km, phút - mét).

%% DANG 9D1?3-2
\textbf{Dạng 2. Toán năng suất, làm chung – làm riêng}
\begin{enumerate}
\item Coi toàn bộ công việc là $1$ đơn vị công việc.
\item Gọi thời gian làm một mình xong công việc của đội 1 là $x$ (giờ/ngày), đội 2 là $y$ ($x, y > 0$).
\item Năng suất trong $1$ đơn vị thời gian: Đội 1 làm được $\dfrac{1}{x}$ (công việc), Đội 2 làm được $\dfrac{1}{y}$ (công việc).
\item Hai đội cùng làm: trong $1$ giờ cả hai làm được $\dfrac{1}{x} + \dfrac{1}{y}$ (công việc).
\item Lập hệ phương trình dạng đặt ẩn phụ $u = \dfrac{1}{x}, v = \dfrac{1}{y}$.
\end{enumerate}

%% DANG 9D1?3-3
\textbf{Dạng 3. Toán về số, chữ số, tỉ lệ, phần trăm}
\begin{enumerate}
\item Số tự nhiên có hai chữ số: $\overline{ab} = 10a + b$ ($a, b \in \mathbb{N}, 1 \le a \le 9, 0 \le b \le 9$). Khi đổi chỗ hai chữ số: $\overline{ba} = 10b + a$.
\item Tăng giảm phần trăm:
\begin{itemize}
\item Giá trị sau khi tăng $r\%$: $x(1 + r\%)$.
\item Giá trị sau khi giảm $r\%$: $x(1 - r\%)$.
\end{itemize}
\end{enumerate}

%% DANG 9D1?3-4
\textbf{Dạng 4. Toán thực tế (hóa học, vật lí, kinh tế)}
\begin{enumerate}
\item \textbf{Bài toán nồng độ dung dịch}:
$$\text{Khối lượng chất tan} = \text{Khối lượng dung dịch} \times \text{Nồng độ \%} \implies m_{ct} = m_{dd} \cdot C\%.$$
\item \textbf{Bài toán kinh tế (giá bán, khuyến mãi, mua sắm)}:
$$\text{Tổng tiền} = \text{Số lượng mặt hàng A} \times \text{Đơn giá A} + \text{Số lượng B} \times \text{Đơn giá B}.$$
\end{enumerate}
